\section{The classical upper bound}\label{sec:upper}

We first prove the upper bound, independently of the construction used
for the lower bound.  The estimate is classical
\cite{AKS1980}.  Our proof uses the uniform-independent-set method of
Alon~\cite{Alon1996}; see also Shearer~\cite{Shearer1983} for the
triangle-free independence bound and Davies, Jenssen, Perkins, and
Roberts~\cite{DJPR2018} for sharp occupancy estimates.

\begin{theorem}\label{thm:upper}
There is an absolute constant $C$ such that, for every integer $t\ge2$,
\[
 r(5,t)\le C\frac{t^4}{(\log t)^3}.
\]
\end{theorem}

The proof has two ingredients.  A triangle-free graph has an independent
set larger than the greedy bound by a logarithmic factor.  A graph whose
adjacent pairs have few common neighbors contains a sufficiently large
triangle-free subgraph after random sampling.  Applying these facts
successively to $K_4$-free and $K_5$-free graphs gives the three powers of
$\log t$.

\begin{lemma}\label{lem:occupancy}
Let $H$ be a finite triangle-free graph on $n$ vertices, and let $D\ge1$
be a real upper bound on its maximum degree.  Then
\[
 \alpha(H)\ge\frac{n\log(D+1)}{8(D+1)}.
\]
\end{lemma}

\begin{proof}
We may assume $n>0$.  Choose an independent set $I$ uniformly from all
independent sets of $H$, including the empty set, and put
$\mu=\mathbb E|I|/n$.  Fix a vertex $x$ and condition on $I$ outside the
closed neighborhood $N[x]=N(x)\cup\{x\}$.  Let $Z_x$ be the number of
neighbors of $x$ having no neighbor in this fixed outside set.  These
$Z_x$ vertices are pairwise nonadjacent.  The possible completions inside
$N[x]$ are therefore $\{x\}$ and the $2^{Z_x}$ subsets of those neighbors.
It follows that
\[
 \Pr(x\in I\mid I\setminus N[x])
 =\frac1{1+2^{Z_x}}\ge\frac12\,2^{-Z_x},
 \qquad
 \mathbb E(|I\cap N(x)|\mid I\setminus N[x])\ge\frac{Z_x}{4}.
\]
The second inequality follows by summing over the subsets; its left side
is $Z_x2^{Z_x-1}/(1+2^{Z_x})$, with value zero when $Z_x=0$.

Now choose $X$ uniformly from the vertices, independently of $I$.
Double counting gives
\[
 \mathbb E Z_X
 \le\frac4n\sum_v\deg(v)\Pr(v\in I)\le4D\mu.
\]
Jensen's inequality for the convex function $z\mapsto2^{-z}$ yields
\begin{equation}\label{eq:occupancy-fixed-point}
 \mu\ge\frac12\mathbb E2^{-Z_X}
 \ge\frac12\,2^{-4D\mu}.
\end{equation}
Write $u=D+1$.  If $\mu<\log u/(8u)$, then
$4D\mu\log2<\tfrac12\log u$, so
\eqref{eq:occupancy-fixed-point} gives $\mu>1/(2\sqrt u)$.
But $\log u\le\sqrt u$ for $u\ge1$, and the assumed inequality gives
$\mu<1/(8\sqrt u)$, a contradiction.
Finally, $\alpha(H)\ge\mathbb E|I|=n\mu$.
\end{proof}

The degree parameter in Lemma~\ref{lem:occupancy} need not equal the
actual maximum degree.  In particular, the lemma also applies to
edgeless graphs by taking any $D\ge1$.

\begin{lemma}\label{lem:sample-triangle-free}
Let $G$ be a finite graph on $n$ vertices with maximum degree $d>0$.
Suppose every adjacent pair has at most $m$ common neighbors, where
$m>0$, $dm\ge1$, and $p=(dm)^{-1/2}$ satisfies $pd\ge1$.
Then $G$ has an induced triangle-free subgraph $H$ with
\[
 |V(H)|\ge\frac34pn,
 \qquad \Delta(H)\le12pd.
\]
\end{lemma}

\begin{proof}
Write $e(G)$ and $T(G)$ for the numbers of edges and triangles.
Each triangle is counted at its three edges, so
\[
 3T(G)=\sum_{xy\in E(G)}|N(x)\cap N(y)|
 \le me(G)\le\frac{ndm}{2}.
\]
Retain each vertex independently with probability $p$, obtaining a set
$S$.  By linearity of expectation,
\begin{align*}
 \mathbb E\left[|S|-T(G[S])-\frac{2e(G[S])}{12pd}\right]
 &\ge pn-\frac{p^3ndm}{6}-\frac{p^2nd}{12pd}\\
 &=\frac34pn.
\end{align*}
Choose a realization at least as large as this expectation.  Delete all
vertices having degree greater than $12pd$ in the original graph
$G[S]$; there are at most $2e(G[S])/(12pd)$ such vertices.  Then delete
one vertex from each remaining triangle until no triangle remains.
The latter procedure uses at most $T(G[S])$ further vertices.
The remaining induced graph has the required order and maximum degree.
\end{proof}

\begin{proof}[Proof of Theorem~\ref{thm:upper}]
We prove the estimates for sufficiently large $t$; increasing the final
constant covers the finitely many smaller values.
First, if $G$ is triangle-free and $\alpha(G)<t$, every neighborhood is
independent, so its maximum degree is at most $t-1$.
Lemma~\ref{lem:occupancy}, with $D=t-1$, gives
$|V(G)|<8t^2/\log t$.  Consequently
\begin{equation}\label{eq:ramsey-three-upper}
 r(3,t)\le C_3\frac{t^2}{\log t}
\end{equation}
for an absolute constant $C_3$.

Next let $s\in\{4,5\}$, and suppose $G$ is $K_s$-free with
$\alpha(G)<t$.  Put $n=|V(G)|$ and $d=\Delta(G)$.
Every neighborhood is $K_{s-1}$-free and has independence number below
$t$, whence $d<r(s-1,t)$.  The common neighborhood of an edge is
$K_{s-2}$-free.  Thus an upper bound on every adjacent codegree is
\[
 m=\begin{cases}
 t,&s=4,\\
 r(3,t),&s=5.
 \end{cases}
\]
Here the first case uses the fact that this common neighborhood is
independent, and the second uses that it is triangle-free.

If $d\le t^{s-5/2}$, greedy selection gives $n<t(d+1)$: repeatedly choose
a vertex for an independent set and remove its closed neighborhood,
which has at most $d+1$ vertices.  Hence
\[
 n\le2t^{s-3/2}
 =o\!\left(\frac{t^{s-1}}{(\log t)^{s-2}}\right).
\]
Otherwise set $p=(dm)^{-1/2}$.  For $s=4$ we have $pd>t^{1/4}$.
For $s=5$, \eqref{eq:ramsey-three-upper} gives
\[
 pd=\sqrt{d/m}
 \ge C_3^{-1/2}t^{1/4}\sqrt{\log t}\ge t^{1/5}
\]
for all sufficiently large $t$.
The hypotheses of Lemma~\ref{lem:sample-triangle-free} therefore hold.
Its subgraph $H$ still satisfies $\alpha(H)<t$.
Applying Lemma~\ref{lem:occupancy} with $D=12pd$ gives
\[
 t>\alpha(H)
 \ge\frac{(3/4)pn\log(12pd+1)}{8(12pd+1)}
 \ge c\frac{n\log t}{d}
\]
for an absolute $c>0$.  We conclude that
\begin{equation}\label{eq:ramsey-upper-step}
 n\le C\frac{td}{\log t}
 \le C\frac{t\,r(s-1,t)}{\log t}
\end{equation}
in the large-degree case.

Apply this argument first with $s=4$, using
\eqref{eq:ramsey-three-upper}.  Together with the small-degree case it
gives $r(4,t)\le C_4t^3/(\log t)^2$.
Apply it next with $s=5$, using this established bound for $r(4,t)$ in
\eqref{eq:ramsey-upper-step}.  This gives
$r(5,t)\le Ct^4/(\log t)^3$ as required.
In each passage from an avoiding graph to a Ramsey number, the
additive one is absorbed into the absolute constant.
\end{proof}
